\subsection{Кодирование, сжатие и пропуск данных}

Производительность колоночного формата определяется тремя группами механизмов:
кодированием значений (компактное двоичное представление, использующее структуру
данных), пропуском нерелевантных данных (отказ от чтения и декодирования блоков,
заведомо не нужных запросу) и векторизованным декодированием (распаковка плотных
массивов значений с использованием параллелизма процессора). Рассмотрим их.

Базовые кодеки. Словарное кодирование (dictionary encoding): для
столбца с небольшим числом различных значений строится словарь, а сами значения
заменяются индексами в нём; выгодно для категориальных строк (страна, тип
устройства, метод запроса). Кодирование длин серий (run-length encoding,
RLE): последовательность одинаковых подряд идущих значений заменяется парой
«значение, длина серии»; выгодно для отсортированных или малоизменчивых столбцов.
Упаковка в биты (bit-packing): целые числа из ограниченного диапазона
хранятся в минимально необходимом числе бит, а не в полном машинном слове.
Кодирование относительно опорного значения (frame of reference, FOR): из
всех значений вычитается минимум блока, остатки бит-упаковываются~\cite{zukowski2006superscalar, lemire2015decoding}. Дельта-кодирование:
хранятся разности соседних значений; выгодно для монотонных последовательностей
(идентификаторы, временные метки). Поверх кодеков применяется блочное сжатие общего
назначения (Snappy, ZSTD, LZ4). Parquet и ORC используют этот арсенал; различие в
деталях (размеры блоков, выбор кодека, наличие фильтров Блума).

FastLanes. Vortex применяет для целочисленных данных схему бит-упаковки
FastLanes~\cite{afroozeh2023fastlanes}, спроектированную под SIMD-декодирование. Обычная бит-упаковка требует при
распаковке битовых сдвигов и масок с зависимостями по данным, что плохо ложится на
векторные инструкции. FastLanes переставляет биты значений в памяти («транспонирует»
их по полосам фиксированной ширины) так, чтобы распаковка \(k\) значений выполнялась
небольшим числом векторных операций без ветвлений и без зависимостей между полосами.
В сочетании с FOR и дельта-кодированием это даёт высокую пропускную способность
декодирования целочисленных столбцов.

ALP. Для чисел с плавающей точкой Vortex применяет кодек ALP~\cite{afroozeh2024alp} (Adaptive
Lossless floating-Point). Идея: многие вещественные значения в реальных данных ---
это десятичные дроби с небольшим числом знаков (координаты, цены, измерения),
которые при умножении на подходящую степень десяти становятся целыми. ALP подбирает
по выборке блока показатели степени, представляет основную массу значений как целые
числа (далее --- FOR и бит-упаковка), а немногочисленные значения, не укладывающиеся
в эту модель (исключения), хранит отдельно в исходном виде. Преобразование
обратимо без потерь. На столбцах с «настоящими» произвольными double ALP
вырождается, но на типичных прикладных данных существенно выигрывает у побайтового
хранения.

FSST. Для строк Vortex применяет кодек FSST~\cite{boncz2020fsst} (Fast Static Symbol Table). По
выборке строит таблицу из не более чем \num{255} коротких подстрок («символов»,
длиной до 8 байт), часто встречающихся в данных, и кодирует строки
последовательностью однобайтовых ссылок на символы (байт \texttt{0xFF} экранирует
буквальный байт). Декодирование --- простая подстановка по таблице. FSST особенно
эффективен на строках с регулярной структурой и общими префиксами: URL и пути
(\texttt{/api/v2/cat\_017}), идентификаторы (\texttt{sess\_0000001234},
\texttt{req\_0000000042}), доменные имена, --- то есть на данных, которые в
аналитических таблицах встречаются повсеместно. Существенно: эффективность FSST (и
прочих кодеков) зависит от порядка строк внутри блока. В отсортированном
блоке соседние строки похожи, символы переиспользуются плотно; в перемешанном блоке
выигрыш меньше. Это объясняет, почему в LSM-дереве выгоды Vortex полнее проявляются
на уплотнённых (отсортированных по ключу) данных, чем на свежих файлах \(L_0\).

Каскадирование кодеков. В Vortex кодеки --- композируемые: к уже
закодированному массиву применяется следующий кодек (например, словарь, затем
бит-упаковка индексов словаря, затем FOR). Выбор каскада адаптивен --- по выборке
оценивается, какая комбинация даёт наилучшее сжатие при приемлемой стоимости
декодирования.

Пропуск данных по статистикам. Файловый формат хранит для каждого блока
(группа строк / страница в Parquet, полоса / группа в ORC, блок в Vortex) статистики
--- минимум, максимум, число null. Если предикат запроса (\texttt{col BETWEEN a AND
b}, \texttt{col = v}, \texttt{col IS NOT NULL}) несовместим с диапазоном блока, блок
не читается. Та же информация на уровне всего файла дублируется в манифесте Paimon,
что позволяет отсечь файл ещё до его открытия. Тонкость: статистики по строковым
столбцам часто хранятся усечёнными (truncate, например, до 16 байт первого
несовпадающего символа) ради компактности; для столбцов, у которых различающая
часть выходит за границу усечения, пропуск перестаёт работать. В работе используется
режим полных статистик (\texttt{metadata.stats-mode = full}), чтобы исключить этот
эффект из сравнения.

Проталкивание предикатов. Помимо пропуска целых блоков, формат может
проталкивать предикат в декодер: вычислять условие на закодированных или
частично декодированных данных и выдавать только подходящие строки, не материализуя
остальные. Vortex поддерживает проталкивание предикатов на уровне нативной
библиотеки: Paimon-предикат преобразуется конвертером в нативное выражение Vortex
(операции сравнения \(=, \ne, <, \le, >, \ge\), проверка на null, конъюнкция и
дизъюнкция для всех поддерживаемых типов), и фильтрация выполняется в Rust-коде до
пересечения границы JNI. Это даёт выигрыш на запросах с селективным предикатом:
через границу JVM/native передаётся лишь малая доля строк.

Отсутствие проталкивания агрегатов. Принципиальное ограничение: интерфейс
файловых форматов Apache Paimon (\texttt{FileFormat}) не предусматривает
проталкивание агрегатов --- операции \texttt{SUM}, \texttt{COUNT}, \texttt{COUNT
DISTINCT}, \texttt{GROUP BY} выполняются вычислительным движком (Flink) над уже
декодированными данными в JVM. Это архитектурное ограничение самого Paimon, а не
конкретного формата. Для Parquet, у которого весь путь --- от декодера до движка
агрегации --- находится в JVM, граница JNI отсутствует. Для Vortex данные сначала
пересекают границу JNI (хотя и без копирования, через Arrow C Data Interface), а
затем агрегируются в JVM; на запросах вида «полное сканирование + агрегация без
селективного предиката» это даёт Vortex дополнительные накладные расходы. Устранение
ограничения --- проталкивание агрегатов в нативный SIMD-слой Vortex --- лежит за
рамками работы и обозначено как направление развития Apache Paimon.

Векторизованное декодирование. И Parquet, и ORC, и Vortex выдают данные
движку не построчно, а пакетами (vectorized batches) --- плотными массивами значений
столбцов, к которым движок применяет векторизованные операции. Vortex дополнительно
строит декодированные пакеты в раскладке Apache Arrow~\cite{arrow-cdata}, что устраняет промежуточное
копирование при передаче в Arrow-совместимые компоненты. Однако в Paimon результат
всё равно адаптируется к внутреннему представлению строк Flink, поэтому полностью
исключить накладные расходы границы не удаётся.
